% The article for the paper-build pull request. Self-contained.
\documentclass[11pt]{article}
\usepackage[T1]{fontenc}
\usepackage{lmodern}
\usepackage{microtype}
\usepackage[margin=1in]{geometry}
\usepackage{amsmath,amssymb,amsthm,mathtools}
\usepackage{algorithm}
\usepackage{algpseudocode}
\usepackage{booktabs}
\usepackage{alltt}
\usepackage[numbers,sort&compress]{natbib}
\usepackage[hidelinks]{hyperref}
\usepackage[capitalise,noabbrev]{cleveref}

\theoremstyle{plain}
\newtheorem{proposition}{Proposition}[section]
\theoremstyle{definition}
\newtheorem{definition}[proposition]{Definition}

\newcommand{\code}[1]{\texttt{#1}}
\newcommand{\pass}{\(\checkmark\)}
\newcommand{\fail}{\(\times\)}
\newcommand{\concat}{\mathbin{+\!\!+}}

\title{Failure Reports and Output Paths for the Paper Build}
\author{Alberto Serrano-Calva}
\date{\today}

\begin{document}
\maketitle

\begin{abstract}
Every pull request in mind-palace carries a LaTeX article, and a continuous-integration job builds
all of them with a shell script around latexmk. Two defects in that script hid failures. When
latexmk failed, the script printed only the path of the TeX log, a file that stays on the runner
and is discarded with it, so a failure that only the runner produced arrived with no reason. An
article that set T1 encoding without Latin Modern failed this way: the runner's TeX Live has no
cm-super, so the EC fonts fall back to bitmaps, and microtype's font expansion stops on a bitmap
font, while a full local TeX Live built the article. Second, an absolute output directory was
joined to the repository root, so latexmk wrote under the repository while the check for undefined
references read a log that did not exist, and the read error counted as a pass. We resolve the
output directory once, make a missing log a failure, and on failure print latexmk's error summary,
the TeX log's error lines with their context and the log's last 30 lines, keeping that log for the
workflow to upload. Six tests run the script on a broken paper, a paper with an undefined
citation, the template with an absolute and a relative output directory, and a stand-in latexmk
that writes no log; five of them fail against the previous script. A local reproduction of the
runner's missing fonts shows the previous report naming a path and the new one naming the pdfTeX
error.
\end{abstract}

\section{Introduction}
Every pull request in mind-palace carries a short article under \code{papers/}, and a
continuous-integration (CI) workflow builds all of them with \code{scripts/build\_papers.sh}. The
script runs latexmk on each article and then reads the TeX log for undefined references and
citations. Its verdict, one line per paper and an exit status, is what leaves the runner: the
workspace and the logs in it are discarded when the job ends, and the workflow uploaded only the
PDFs.

Two defects in the script hid failures. When latexmk failed, the script printed the path of the
log, a file the reader of the job's output could not open. When the output directory was
absolute, it built the paper under the repository and then checked a log that did not exist, and
the failed read passed the check. The first defect turned a failure that only the runner could
produce into a guess; the second could pass a paper with undefined citations. Contributions:
\begin{itemize}
  \item an account of both defects: a report that names a file on the runner, and write and read
        paths that diverge for an absolute output directory, with a read error that counts as a
        pass (\cref{sec:defects});
  \item a fix that resolves the output directory once, fails a paper that leaves no log, and prints
        the reason from the log that explains the failure, keeping that log for the workflow to
        upload (\cref{sec:fix});
  \item six tests, five of which fail against the previous script, and a reproduction of the
        failure that motivated the change, with a null control (\cref{sec:eval}).
\end{itemize}

\section{Background}
\paragraph{latexmk.} latexmk runs \code{pdflatex} and BibTeX as rules, repeating them until the
cross-references settle; with \code{-outdir} it writes the output, the log included, to one
directory \cite{collins2026latexmk}. When a rule fails, latexmk ends its output with a summary of
the errors it collected, naming the rule: \code{pdflatex} for a TeX error, \code{bibtex} for a
bibliography error.

\paragraph{TeX's log.} TeX reports an error in a fixed form: a line beginning with \code{!} states
it, and the lines after it show the context, the current input line broken where TeX stopped and
prefixed \code{l.}\textit{n} with its line number \cite{knuth1984texbook}. With
\code{-halt-on-error} the first error ends the run, so the error sits near the end of the log,
followed by a summary of TeX's memory use. TeX Live wraps log lines at 79 columns.

\paragraph{Font expansion.} pdfTeX can stretch and shrink glyphs slightly to even out the spacing
of justified lines, following Zapf's hz-program \cite{thanh2000microtypography,pdftexmanual}. In
the recommended \code{autoexpand} mode it scales the glyphs of an outline font within the page
stream \cite{pdftexmanual}, and the microtype package enables expansion by default where it can
assume it works \cite{schlicht2026microtype}. A bitmap font cannot be scaled this way, and pdfTeX
stops with \emph{auto expansion is only possible with scalable fonts}.

\paragraph{Fonts on two machines.} T1 encoding selects the EC fonts unless a package such as
lmodern selects others. The EC fonts exist as Type~1 outlines only through cm-super
\cite{volovich2008cmsuper}; without it, TeX Live generates bitmaps from their METAFONT sources.
Latin Modern supplies Type~1 fonts in T1 encoding itself \cite{jackowski2021lm}. The full TeX Live
scheme includes cm-super. The CI runner installs Ubuntu's TeX Live packages (\code{latexmk},
\code{texlive-latex-recommended}, \code{texlive-latex-extra}, \code{texlive-fonts-recommended},
\code{texlive-science}, \code{lmodern}) without recommended packages, on an Ubuntu 26.04 image with
TeX Live 2025 and latexmk 4.87. The job's log lists the 24 packages it set up, and cm-super is not
among them. An article that sets T1 encoding without Latin Modern therefore gets outline fonts
locally and bitmaps on the runner, and with microtype it builds locally and stops on the runner.
One article did.

\paragraph{Related work.} The reason a CI build failed is usually a short chunk of a long log;
LogChunks annotates 797 build logs with the chunk that explains each failure
\cite{brandt2020logchunks}. Here the log's format is known, so the chunk can be selected by TeX's
own conventions instead of learned. A check that counts an unreadable input as a pass has the
opposite of a fail-safe default \cite{saltzer1975protection}.

\section{The defects}\label{sec:defects}
\paragraph{The report.} The script sent latexmk's output to \code{/dev/null} and, on failure,
printed
\begin{small}
\begin{alltt}
\fail{} 2026-11-07-graph-dynamics: see build/papers/2026-11-07-graph-dynamics/paper.log
\end{alltt}
\end{small}
The path is relative to a workspace that is gone once the job ends, and the workflow did not upload
logs. What else reached the job's output was latexmk's progress and advice on standard error,
ending in six lines that suggest its \code{-f} option, the same for every failure. The reason had to
be inferred from the difference between the two TeX installations.

\paragraph{The output path.} Let $r$ be the absolute path of the repository root, $p$ the value of
\code{PAPERS\_OUT} and $n$ the name of a paper's directory.

\begin{definition}
The \emph{resolution} of $p$ against $r$ is $\rho(r,p) = p$ if $p$ is absolute and $r/p$
otherwise. The \emph{concatenation} $r \concat p$ is the string $r$, a slash, then $p$.
\end{definition}

The script changed to $r$, then, from inside the paper's directory, gave latexmk the output
directory \code{\$OLDPWD/\$OUT/\$name}, where \code{\$OLDPWD} is $r$. latexmk therefore wrote the
log under $W(p) = (r \concat p)/n$. The check read \code{\$OUT/\$name/paper.log} with the working
directory at $r$, that is, under $R(p) = \rho(r,p)/n$.

\begin{proposition}\label{prop:paths}
$W(p) = R(p)$ if and only if $p$ is relative.
\end{proposition}

\begin{proof}
For a relative $p$, $r \concat p = r/p = \rho(r,p)$. For an absolute $p$, $r \concat p$ is $r$
followed by $p$, a path inside $r$ (the doubled slash is one separator), while $\rho(r,p) = p$;
the two differ, since the first is the longer.
\end{proof}

So an absolute $p$ sent the build under the repository, at a path that repeats $p$ inside it, and
left at $R(p)$ only the empty directory the script had made there.

\begin{proposition}\label{prop:open}
With an absolute $p$ and an output directory not written before, the previous check passed every
paper that latexmk built, whatever its log said.
\end{proposition}

\begin{proof}
By \cref{prop:paths} the log is under $W(p)$, and no file exists at $R(p)/\code{paper.log}$.
\code{grep} exits with status 0 when a line is selected, 1 when none is, and greater than 1 on an
error such as a missing file \cite{posix2024}. The script tested that status with \code{if}, which
takes every nonzero status as false, so a missing file and a log free of undefined references led
to the same branch, which printed \pass.
\end{proof}

The check conflated two statements: the log has no undefined references, and the log could not be
read. A check over a build's log is sound only if it reads the log the build wrote and counts a
failure to read it as a failure.

\section{The fix}\label{sec:fix}
The script now resolves the output directory once, $o = \rho(r, p)$, after changing to $r$, and
derives every path from $o$: latexmk's output directory, the log the check reads and the copy kept
on failure. The write and read paths are then equal by construction. A relative $p$ keeps its
meaning, a path from the repository root, as the paper directories given as arguments are.

\begin{algorithm}
\caption{Building one paper directory $d$ into $o/n$.}\label{alg:build}
\begin{algorithmic}[1]
\State remove $o/n/\code{failed.log}$
\State run latexmk on $d$ with output directory $o/n$, its own output to $o/n/\code{latexmk.log}$
\If{latexmk failed}
  \If{latexmk's output holds its error summary \textbf{and} $o/n/\code{paper.log}$ exists}
    \State $L \gets o/n/\code{paper.log}$ \Comment{a rule ran and failed}
    \State print the summary, then the error lines of $L$ with two lines of context each
    \State print the last 30 lines of $L$
  \Else
    \State $L \gets o/n/\code{latexmk.log}$ \Comment{no rule ran; a TeX log there is old}
    \State print the last 30 lines of $L$
  \EndIf
\ElsIf{$o/n/\code{paper.log}$ does not exist}
  \State $L \gets o/n/\code{latexmk.log}$; report the missing log \Comment{fail closed}
\ElsIf{the log reports an undefined reference or citation}
  \State $L \gets o/n/\code{paper.log}$; print each such warning once
\Else
  \State print \pass{} and stop
\EndIf
\State copy $L$ to $o/n/\code{failed.log}$; the paper fails
\end{algorithmic}
\end{algorithm}

An error line is one that starts with \code{!}, contains \code{pdfTeX error}, or starts with a
\textit{file}\code{:}\textit{line}\code{:} prefix, the form TeX uses when asked for file and line
errors. Two choices in \cref{alg:build} need a reason. latexmk's summary decides which log
to show because the TeX log in the output directory may be from an earlier run: if latexmk cannot
start at all, the directory still holds the last successful build's log, and its tail would
describe a success. latexmk prints its summary only after one of its rules ran and failed, and then
the TeX log is the one that run wrote or kept. The two lines after each error line are where TeX
prints the context, the input line and the point at which it stopped; the last 30 lines cover an
error that the line match misses.

latexmk's own output, part of which used to reach the console on standard error, now goes to
\code{latexmk.log}, so a passing build prints one line per paper. The workflow uploads
\path{build/papers/*/paper.pdf} and \path{build/papers/*/failed.log} as one artifact whose root
is \path{build/papers} \cite{uploadartifact}, so each failed paper's log sits in its directory
beside the other papers' PDFs. After the build, the same job runs the script's tests, since it has
TeX Live installed.

\paragraph{The tests.} \path{tests/integration/test_build_papers.py} runs the script in a copy
of the repository's layout under a temporary directory, from a working directory outside the copy,
so a relative output directory resolves inside the copy and nothing is written to the checkout.
Five tests build with TeX and carry the marker \code{needs\_tex}; they skip where latexmk is
absent \cite{pytestdocs}, and the workflow selects them with \code{-m needs\_tex}. The sixth puts a
stand-in latexmk first on the path, one that succeeds without writing a log, so it needs no TeX and
runs in every pytest tier. \Cref{tab:tests} lists them.

\section{Evaluation}\label{sec:eval}
The runs below are on one machine with TeX Live 2026 (pdfTeX 1.40.29, latexmk 4.88, cm-super
installed).

\paragraph{The failure that motivated the change.} The dynamics article failed on the runner and
built locally; the commit that fixed it added \code{\textbackslash usepackage\{lmodern\}}. To give
the local machine the runner's fonts for one build, we put
\code{\textbackslash pdfmapfile\{-cm-super-t1.map\}} at the top of the article. It deletes
cm-super's T1 entries from pdfTeX's font map before any font is used \cite{pdftexmanual}, so the EC
fonts have no outline and TeX Live generates bitmaps, as on the runner. With this line, the article
as it was before its fix, built by the previous script, gave the runner's report exactly:
\begin{small}
\begin{alltt}
\fail{} 2026-11-07-graph-dynamics: see build/papers/2026-11-07-graph-dynamics/paper.log
\end{alltt}
\end{small}
with latexmk's six lines of advice on standard error. The new script printed latexmk's summary
(\code{pdflatex} returned 1), then
\begin{small}
\begin{alltt}
! pdfTeX error (font expansion): auto expansion is only possible with scalable
fonts.
\end{alltt}
\end{small}
followed by TeX's context and the log's last 30 lines, which place the error at the first page's
shipout. As a null control, the article after its fix, with the same line added, built and passed,
and its log names no bitmap font.

\paragraph{The output path.} On the undefined-citation fixture with an absolute output directory
outside the repository's copy, the previous script printed \pass{} and exited 0; \code{grep}
reported on standard error that the log did not exist, and the PDF was written inside the
repository's copy, under a tree that repeats the absolute path. The new script exited 1, printed
natbib's warning naming the missing key, and wrote only under the output directory.

\paragraph{The tests against both scripts.} \Cref{tab:tests} runs the six tests against the
previous script and the new one. Against the previous script five fail. The one that passes builds
the template into a relative output directory, the case \cref{prop:paths} says the previous script
handled, so it serves as the control. The relative undefined-citation test fails for a different
reason than the absolute one: the previous script did fail that paper, but its report named no
key. Against the new script all six pass, in 12.9\,s and 13.8\,s over two runs.

\begin{table}[htbp]
\centering
\caption{The six tests against the previous script and the new one.}\label{tab:tests}
\begin{tabular}{@{}lccc@{}}
\toprule
test & needs TeX & previous & new \\
\midrule
template, absolute \code{PAPERS\_OUT}, writes only there & yes & fail & pass \\
template, relative \code{PAPERS\_OUT}, writes only there & yes & pass & pass \\
undefined citation, absolute & yes & fail & pass \\
undefined citation, relative & yes & fail & pass \\
broken paper: error and context printed, next paper built & yes & fail & pass \\
no log written: the paper fails & no & fail & pass \\
\bottomrule
\end{tabular}
\end{table}

\paragraph{Other failures.} A bibliography with a syntax error makes BibTeX fail. The TeX log
then holds no error line, and the report shows latexmk's summary, which names the \code{bibtex}
rule and its \code{.blg} file. With latexmk missing from the path, the report shows the shell's
\emph{command not found} from \code{latexmk.log}; reading the TeX log instead would have shown the
tail of the last successful build's log.

\paragraph{The repository's papers and gates.} Every article in the repository, the template and
this article built and passed with the new script, both into the default directory and into an
absolute one. The repository's gates passed on the final tree (\cref{tab:gates}); the mypy floors
are the packages \code{core}, \code{agents}, \code{eval}, \code{ops}, \code{scheduler} and
\code{scripts}.

\begin{table}[htbp]
\centering
\caption{The repository's gates on this change. No test failed.}\label{tab:gates}
\begin{tabular}{@{}ll@{}}
\toprule
gate & result \\
\midrule
ruff & clean \\
import firewall & clean \\
mypy over the floors & 0 errors \\
mypy over the whole repository & 69 errors, the pinned baseline \\
type gate & clean \\
deterministic pytest tier & 3{,}420 passed, 11 skipped, 21 deselected \\
\bottomrule
\end{tabular}
\end{table}

\section{Limitations}
The excerpt follows TeX's conventions for errors. A failure that leaves no error line in the TeX
log, such as BibTeX's, is reported through latexmk's summary, which names the \code{.blg} file;
that file is not uploaded. TeX wraps log lines at 79 columns, so an error line can be split, as
the font-expansion error is; the context lines that follow carry the rest. A citation warning
whose key pushes the word \emph{undefined} past the wrap is not matched on its own line, but
LaTeX's closing warning that there were undefined citations still fails the paper. The choice of
log reads latexmk's wording, \emph{Collected error summary}; if a later latexmk changes it,
failures fall to the branch that prints latexmk's output, which still contains pdflatex's
messages. The reproduction deletes one map file's entries for one run; it reproduces the
font-expansion failure, not the runner's TeX Live. And the fix makes the difference between the two
TeX installations visible; it does not remove it.

\section{Future work}
A check over \code{papers/*/paper.tex} that rejects T1 encoding without Latin Modern would catch
the motivating failure before the workflow does, and building locally with the runner's packages,
in a container, would catch the whole class. A workflow annotation carrying the first error line
would put the reason on the pull request's page. A longer TeX line for the build, through
\code{max\_print\_line}, would keep each error and warning on one line. BibTeX's log could be kept
with the TeX log when the \code{bibtex} rule fails.

\section{Conclusion}
The paper build now says why a paper failed, in the job's output and in an uploaded log, and its
check reads the log the build wrote whatever form the output path takes; a missing log fails the
paper. The failure that motivated the change, rebuilt with the runner's fonts simulated, now
reports pdfTeX's error instead of a path.

\bibliographystyle{plainnat}
\bibliography{refs}

\end{document}
